\documentclass[10pt,a4paper]{amsart}
\usepackage[margin=19mm]{geometry}
\usepackage{amsmath,amssymb,mathtools}
\usepackage[scr=rsfs,cal=euler]{mathalfa}
\usepackage{xfrac}
\usepackage[unicode,hidelinks,bookmarks=false]{hyperref}
\newcommand{\ie}{i.e.\ }
\newcommand{\cf}{cf.\ }
\newcommand{\resp}{respectively\ }
\newcommand{\Subsection}{\subsection}
\providecommand{\nobreakdash}{\nobreak\hbox{-}}
\setlength{\emergencystretch}{2em}
\multlinegap0pt
\usepackage{bm}
\usepackage{bbm}
\usepackage{dsfont}
\usepackage{xspace}
\usepackage{cancel}
\usepackage{mathtools}
\usepackage{amsthm}
\makeatletter
\@ifundefined{definition}{\newtheorem{definition}{Definition}}{}
\@ifundefined{proposition}{\newtheorem{proposition}{Proposition}}{}
\@ifundefined{theorem}{\newtheorem{theorem}{Theorem}}{}
\makeatother
\newcommand{\edom}{\gamma_{e}}
\title[Domination versus edge domination in regular graphs of degre]{Domination versus edge domination in regular graphs of degree at least seven}
\author{Chakshu Gupta}
\date{Source: arXiv:2608.22498v1 (CC BY 4.0)}
\begin{document}
\maketitle

\noindent\emph{Source and licence.} Domination versus edge domination in regular graphs of degree at least seven. Version arXiv:2608.22498v1 (CC BY 4.0), distributed under the Creative Commons Attribution 4.0 International licence, as recorded in the arXiv licence metadata for this exact version and shown on the abstract page https://arxiv.org/abs/2608.22498v1. Full rights notice: licenses/arxiv-2608.22498-CC-BY-4.0.txt.\par\smallskip

\noindent\textit{Editorial scope.} Counted unit: the Theorem whose source label is \texttt{thm:lll} in the article (source0.tex lines 327--330, source label \texttt{thm:lll}), delivered together with the article's own complete original proof of it (source0.tex lines 332--400, one \texttt{proof} environment of 69 lines). No mathematics is written, repaired or re-derived: statement and proof are the article's own text line for line. Required context retained uncounted: prop:reduction (lines 185--188); def:atrisk (lines 281--285); eq:lopsi (lines 319--323). Every definition, display and label of those lines is retained unchanged. Omitted from this package and not redistributed: 1--184, 189--280, 286--318, 324--326, 331--331, 401--984. Nothing inside a retained excerpt is omitted or altered: every retained body is delivered exactly as the article's lines are written, so no omission receipt of any kind accompanies this package. Citations: the retained excerpts carry no citation command at all, so no bibliography entry of the article is transcribed and no reference is redistributed. Reference adaptation: none was needed and none was made; every reference inside the delivered unit resolves inside it, so no unresolved reference or citation is produced. Printed slips kept literal: no printed word, symbol, hypothesis or number of the counted statement or of its proof was altered. Layout and class: the article's own class is \texttt{article} with packages including geometry, amsmath, amssymb, amsthm, booktabs, xcolor, hyperref, url; the wrapper is the lane's pinned \texttt{amsart} editorial wrapper at 10pt on A4, which restates the theorem-like environments the retained text needs and adds the article's own macro definitions that the retained text needs (\texttt{\textbackslash edom}), with extra wrapper packages (bm, bbm, dsfont, xspace, cancel, mathtools). The delivered numbering is the unit's own. Assets: no retained span contains a figure environment, a graphics-inclusion command, a PDF-page inclusion, a table or a TikZ picture; no image file is copied into this package, the package declares no asset dependency and no image file is redistributed with it. Licence: version arXiv:2608.22498v1 of this article is distributed under the Creative Commons Attribution 4.0 International licence, as recorded in the arXiv licence metadata for this exact version and shown on the abstract page https://arxiv.org/abs/2608.22498v1. A full rights notice is delivered at licenses/arxiv-2608.22498-CC-BY-4.0.txt. Characters: every retained excerpt is pure ASCII, so the delivered engine, XeLaTeX with its default Latin Modern text font, reports no missing character.
\begin{proposition}[Reduction]\label{prop:reduction}
If a minimum maximal matching $M$ of a graph $G$ admits a dominating
transversal, then $\gamma(G) \le |M| = \edom(G)$.
\end{proposition}

\begin{definition}[At-risk vertex]\label{def:atrisk}
Let $M$ be a maximal matching of a $\Delta$-regular graph $G$ with unsaturated set
$Z$. A vertex $u \in Z$ is at risk if its $\Delta$ neighbours lie on $\Delta$
distinct $M$-edges.
\end{definition}

\begin{equation}\label{eq:lopsi}
  \Pr\Bigl[A_i \,\Big|\, \textstyle\bigwedge_{j \in S} \overline{A_j}\Bigr]
  \;\le\; \Pr[A_i]
  \qquad\text{for every } i \text{ and every } S \subseteq V(H)\setminus N_H[i].
\end{equation}

\begin{theorem}[Local Lemma bound]\label{thm:lll}
Every $\Delta$-regular graph $G$ with $\Delta \ge 7$ satisfies
$\gamma(G) \le \edom(G)$.
\end{theorem}

\begin{proof}
A symmetric application of the Local Lemma
reaches $\Delta \ge 9$, and the lopsided form, which keeps only the conflicting
dependencies and discards the agreeing ones, reaches $\Delta \ge 7$.

\medskip\noindent
\emph{Setup.} Let $M$ be a minimum maximal matching of $G$, so $|M| = \edom(G)$. Choose one
endpoint of each $M$-edge independently and uniformly at random; this is a
random transversal $T$, encoded by independent variables $X_e$, one per $M$-edge,
each uniform on the two endpoints of $e$. For an at-risk vertex $u$, let $A_u$ be
the event that $u$ is uncovered, that is, that none of its $\Delta$ neighbours is
chosen. These neighbours lie on $\Delta$ distinct $M$-edges
(Definition~\ref{def:atrisk}), so $A_u$ is an
intersection of $\Delta$ independent events of probability $\tfrac12$, giving
\[
  \Pr[A_u] = 2^{-\Delta}.
\]
If no event $A_u$ occurs, then $T$ covers every at-risk vertex, hence is
dominating, and Proposition~\ref{prop:reduction} gives
$\gamma(G) \le |M| = \edom(G)$. It remains to show
$\Pr[\bigwedge_u \overline{A_u}] > 0$.

\medskip\noindent
Each $A_u$ is an atom of the $\Delta$ variables $\{X_e : e \text{ incident to a
neighbour of } u\}$, fixing each to leave that neighbour unchosen. Two events
$A_u, A_v$ share a variable $X_e$ exactly when each of $u$ and $v$ has a neighbour
among the endpoints of $e$.

\medskip\noindent
\emph{Symmetric bound.} Every shared-variable dependency is counted. The event
$A_u$ is mutually independent of all $A_v$ sharing no variable with it. For a fixed $M$-edge
$e = \{x, x'\}$ with $u \sim x$, an at-risk $v \neq u$
sharing $X_e$ is a neighbour of $x$ or of $x'$ and, being unsaturated, is
neither endpoint. The
neighbours of $x$ besides $u$ and $x'$ number $\Delta-2$, and those of $x'$
besides $x$ number $\Delta-1$, so at most $2\Delta-3$ per edge and
$\Delta(2\Delta-3)$ in total. The
symmetric Local Lemma applies when $e\,2^{-\Delta}\bigl(\Delta(2\Delta-3)+1\bigr)
\le 1$, which holds for $\Delta \ge 9$, the left side being $\approx 0.72$ at
$\Delta = 9$ and $\approx 1.11$ at $\Delta = 8$.

\medskip\noindent
\emph{Lopsided bound.} Counting only the conflicting dependencies improves the
threshold. Let $H$ join
$A_u \sim A_v$ exactly when they conflict, sharing a variable $X_e$ on which they
demand different endpoints. This conflict graph is a lopsidependency graph by the
standard argument for atomic events. Fix a set $S$ of events none joined to
$A_u$; since $\Pr[A_u] > 0$, condition~\eqref{eq:lopsi} is equivalent to
$\Pr[B \mid A_u] \le \Pr[B]$, where $B = \bigwedge_{v \in S}\overline{A_v}$. Each
such $A_v$ agrees with $A_u$ on every shared variable, so conditioning on
$A_u$ leaves $A_v$ as a weaker event $A_v' \supseteq A_v$ on the remaining variables,
independent of $A_u$; equivalently, $\overline{A_v'} \subseteq \overline{A_v}$. Hence
\[
  \Pr[B \mid A_u] = \Pr\Bigl[\textstyle\bigwedge_{v\in S}\overline{A_v'}\Bigr]
  \le \Pr[B].
\]

\medskip\noindent
In $H$, $A_u$ is joined, for each $M$-edge $e = \{x,x'\}$ with
$u \sim x$, to the at-risk vertices adjacent to the opposite endpoint $x'$, which
demand $X_e = x$ where $A_u$ demands $x'$. The partner $x$ is saturated, hence
not at-risk, so the joined events lie among the other $\Delta - 1$ neighbours of
$x'$, at most $\Delta - 1$ per edge, and $H$
has maximum degree at most $\Delta(\Delta-1)$. The lopsided
Local Lemma applies when $e\,2^{-\Delta}\bigl(\Delta(\Delta-1)+1\bigr) \le 1$,
which holds for $\Delta \ge 7$, the left side being $\approx 0.91$ at
$\Delta = 7$ and $\approx 1.32$ at $\Delta = 6$. Hence
$\Pr[\bigwedge_u \overline{A_u}] > 0$, completing the proof.
\end{proof}

\end{document}
